% One-time imported draft. Editable native source; no Markdown regeneration.
\section{Hexadecimal preserves the binary
  grouping}\label{004-bit-inspector-p01-hexadecimal-preserves-the-binary-grouping}

A binary digit represents one bit. A hexadecimal digit represents four
bits, with digit values zero through fifteen. Eight hexadecimal digits
therefore specify the same 32-bit word as 32 binary digits. Changing the
notation changes the written form, not the stored pattern.

For example, 0x80000001 has the top word bit and the bottom word bit
set. Its binary form is a one, thirty zeroes, and a final one. Its
unsigned value is 2147483649, and its signed value is \ensuremath{-}2147483647. Both
integer interpretations refer to the same hexadecimal pattern.

Leading zeroes make a chosen width visible. The value one can be written
as 0x00000001 and as a 32-digit binary string ending in one. The report
deliberately keeps that width, so a reader can locate bit 31 without
inferring how many leading positions were omitted.
